\documentclass[12pt,a4paper]{article}

\usepackage[T1]{fontenc}
\usepackage[utf8]{inputenc}
\usepackage[romanian]{babel}
\usepackage{geometry}

\geometry{margin=2.5cm}

\title{Tutoriat 4\\Recapitulare si Exercitii}
\author{}
\date{}

\begin{document}

\maketitle



\section{Recapitulare}

În acest tutoriat au fost prezentate conceptele de Overloading, Overriding și Upcasting, împreună cu diferențele și aplicațiile lor în Programarea Orientată pe Obiecte.

\subsection{Overloading}

Overloading (supraîncărcarea):

\begin{itemize}
\item apare în aceeași clasă;
\item funcțiile au același nume;
\item funcțiile au parametri diferiți;
\item nu necesită moștenire.
\end{itemize}

Exemplu:

\begin{verbatim}
void f();

void f(int);

void f(double);
\end{verbatim}

\subsection{Hiding}

Hiding apare atunci când o clasă derivată declară o metodă cu același nume ca o metodă din clasa de bază.

În această situație, metodele cu același nume din baza devin inaccesibile direct.

\begin{verbatim}
class B {
public:
    void f();
};

class D : public B {
public:
    void f(int);
};
\end{verbatim}

\textbf{Observație:} La Hiding contează doar numele metodei. Chiar dacă parametrii diferă, metodele din clasa de bază sunt ascunse.

\subsection{Overriding}

Overriding (suprascrierea):

\begin{itemize}
\item apare la moștenire;
\item metodele au același nume;
\item metodele au aceiași parametri;
\item metoda din derivată înlocuiește metoda din bază.
\end{itemize}

Exemplu:

\begin{verbatim}
class B {
public:
    void f();
};

class D : public B {
public:
    void f();
};
\end{verbatim}

\subsection{Operatorul de rezoluție ::}

Operatorul ::

permite accesarea explicită a membrilor unei clase.

Exemplu:

\begin{verbatim}
d.B::f();
\end{verbatim}

Acesta este utilizat pentru:

\begin{itemize}
\item accesarea metodelor ascunse;
\item accesarea metodelor suprascrise;
\item eliminarea ambiguităților.
\end{itemize}

\subsection{Upcasting}

Upcasting reprezintă conversia unui pointer sau a unei referințe către o clasă derivată într-un pointer sau o referință către clasa de bază.

Exemplu:

\begin{verbatim}
Circle c;

Shape* p = &c;
\end{verbatim}

Caracteristici:

\begin{itemize}
\item conversie automată;
\item conversie sigură;
\item funcționează doar pentru pointeri și referințe.
\end{itemize}

\subsection{Limitările Upcasting-ului}

După Upcasting, accesul este limitat la membrii clasei de bază.

\begin{verbatim}
Shape* p = &c;

p->draw();      // OK
p->setRadius(); // Eroare
\end{verbatim}

\textbf{Observație:} Obiectul își păstrează tipul real, însă metodele accesibile sunt determinate de tipul pointerului sau al referinței.

\subsection{Idei de reținut}

\begin{enumerate}
\item Overloading = aceeași clasă + parametri diferiți.
\item Overriding = moștenire + aceiași parametri.
\item Hiding = moștenire + același nume, dar parametri diferiți.
\item Tipul returnat nu participă la Overloading.
\item Operatorul :: permite accesarea explicită a membrilor unei clase.
\item Upcasting convertește Derivată → Bază.
\item După Upcasting avem acces doar la membrii clasei de bază.
\item Upcasting funcționează doar cu pointeri și referințe.
\end{enumerate}


\section{Exerciții}

\subsection{Exercițiul 1 - Overloading}

Care dintre următoarele declarații reprezintă exemple valide de supraîncărcare?

\begin{verbatim}
void f(int x);

void f(double x);

void f(int x, int y);
\end{verbatim}

Justificați răspunsul.

\subsection{Exercițiul 2 - Tipul returnat}

Analizați următorul cod:

\begin{verbatim}
int f(int x);

double f(int x);
\end{verbatim}

Compilarea are succes?

Motivați răspunsul.

\subsection{Exercițiul 3 - Hiding}

Determinați ce instrucțiuni sunt valide.

\begin{verbatim}
class B {
public:
    void f() {
    }
};

class D : public B {
public:
    void f(int x) {
    }
};

D d;
\end{verbatim}

Apeluri:

\begin{verbatim}
d.f();

d.f(10);
\end{verbatim}

Precizați dacă apar erori de compilare.

\subsection{Exercițiul 4 - Operatorul ::}

Pornind de la exercițiul anterior, scrieți instrucțiunea necesară pentru apelarea metodei:

\begin{verbatim}
B::f()
\end{verbatim}

prin intermediul obiectului:

\begin{verbatim}
d
\end{verbatim}

\subsection{Exercițiul 5 - Overriding}

Analizați următorul cod:

\begin{verbatim}
class B {
public:
    void f() {
        std::cout << "B";
    }
};

class D : public B {
public:
    void f() {
        std::cout << "D";
    }
};
\end{verbatim}

Determinați afișarea pentru:

\begin{verbatim}
D d;

d.f();
\end{verbatim}

\subsection{Exercițiul 6}

Pentru fiecare situație indicați conceptul corespunzător:

\begin{enumerate}
\item Aceeași clasă, aceeași metodă, parametri diferiți.
\item Clasă derivată, aceeași metodă, aceiași parametri.
\item Clasă derivată, aceeași metodă, parametri diferiți.
\end{enumerate}

\subsection{Exercițiul 7}

Considerăm:

\begin{verbatim}
class Shape {
};

class Circle : public Shape {
};

Circle c;

Shape* p = &c;
\end{verbatim}

Identificați:

\begin{itemize}
\item tipul obiectului;
\item tipul pointerului;
\item tipul conversiei realizate.
\end{itemize}

\subsection{Exercițiul 8 - Acces după Upcasting}

Analizați codul:

\begin{verbatim}
class Shape {
public:
    void draw() {
    }
};

class Circle : public Shape {
public:
    void setRadius() {
    }
};

Circle c;

Shape* p = &c;
\end{verbatim}

Indicați care dintre următoarele apeluri sunt valide:

\begin{verbatim}
p->draw();

p->setRadius();
\end{verbatim}

Motivați răspunsul.

\subsection{Exercițiul 9 - Adevărat sau Fals}

Indicați dacă afirmațiile sunt adevărate sau false.

\begin{enumerate}
\item Overloading necesită moștenire.
\item Overriding necesită moștenire.
\item Tipul returnat participă la Overloading.
\item Upcasting este realizat automat.
\item Upcasting funcționează doar cu pointeri și referințe.
\item Hiding poate apărea chiar dacă parametrii diferă.
\end{enumerate}

\subsection{Exercițiul 10 - Recapitulare}

Completați spațiile libere:

\begin{enumerate}
\item Overloading apare în aceeași ............ .
\item Overriding apare la ............ .
\item Hiding ascunde metodele din clasa ............ .
\item Operatorul utilizat pentru acces explicit este ............ .
\item Upcasting realizează conversia de la clasa ............ la clasa ............ .
\item După Upcasting avem acces doar la membrii clasei ............ .
\end{enumerate}


\subsection{Exercițiul 6 - Răspunsuri}
Răspunsurile trebuie să fie:

\begin{itemize}
\item Overloading;
\item Overriding;
\item Hiding.
\end{itemize}

\subsection{Răspunsuri Exercițiul 9}

\begin{enumerate}
\item Fals
\item Adevărat
\item Fals
\item Adevărat
\item Adevărat
\item Adevărat
\end{enumerate}

\subsection{Răspunsuri Exercițiul 10}

\begin{enumerate}
\item Overloading apare în aceeași \textbf{clasă}.
\item Overriding apare la \textbf{moștenire}.
\item Hiding ascunde metodele din clasa \textbf{de bază}.
\item Operatorul utilizat pentru acces explicit este \textbf{::}.
\item Upcasting realizează conversia de la clasa \textbf{derivată} la clasa \textbf{de bază}.
\item După Upcasting avem acces doar la membrii clasei \textbf{de bază}.
\end{enumerate}



\end{document}